\documentclass[11pt]{article}

% Shared notes settings for Elements of AIML lectures (UPES).
% Keep this file free of \documentclass to allow reuse via \input.

\usepackage[utf8]{inputenc}
\usepackage[T1]{fontenc}
\usepackage[a4paper,margin=1in]{geometry}
\usepackage{hyperref}
\usepackage{xurl}
\usepackage{booktabs}
\usepackage{amsmath}
\usepackage{amssymb}
\usepackage{listings}
\usepackage{xcolor}
\usepackage{enumitem}
\usepackage{parskip}

% Code listings (general-purpose).
\lstset{
  basicstyle=\ttfamily\small,
  keywordstyle=\color{blue},
  commentstyle=\color{gray},
  stringstyle=\color{teal},
  showstringspaces=false,
  columns=fullflexible,
  frame=single,
  framerule=0.2pt,
  breaklines=true
}

\setlist[itemize]{topsep=0.3em,itemsep=0.2em}
\setlist[enumerate]{topsep=0.3em,itemsep=0.2em}

% Simple per-section source line.
\newcommand{\notesources}[1]{\par\noindent{\small\color{gray}\textit{Sources:} \sloppy #1\par}}

% Unit-wise source strings for this subject (used by slides/notes).
% Path is relative to the lecture `latex/` folder (the compilation CWD).
\IfFileExists{../../../_shared/aiml-sources.tex}{%
  % Source strings for Elements of AIML (used by slides + notes).

\newcommand{\AIMLRepoURL}{https://github.com/tali7c/Elements-of-AIML}

\newcommand{\AIMLLabSources}{%
Provost and Fawcett, \textit{Data Science for Business}.;%
McKinney, \textit{Python for Data Analysis}.;%
WEKA Documentation (Waikato Environment for Knowledge Analysis), accessed 2026-02-28.;%
Matplotlib and Seaborn documentation, accessed 2026-02-28.%
}

\newcommand{\AIMLUnitISources}{%
Rich and Knight, \textit{Artificial Intelligence}, McGraw Hill, 2017.;%
Russell and Norvig, \textit{Artificial Intelligence: A Modern Approach} (4e), Ch. 1--2 (Introduction, Intelligent Agents).;%
Poole and Mackworth, \textit{Artificial Intelligence: Foundations of Computational Agents} (2e), selected sections.%
}

\newcommand{\AIMLUnitIISources}{%
Russell and Norvig, \textit{Artificial Intelligence: A Modern Approach} (4e), Ch. 7--9 (Logical Agents, First-Order Logic, Inference).;%
Huth and Ryan, \textit{Logic in Computer Science} (2e), selected sections on propositional and first-order logic.;%
Genesereth and Nilsson, \textit{Logical Foundations of Artificial Intelligence}, selected chapters.%
}

\newcommand{\AIMLUnitIIISources}{%
Mueller and Massaron, \textit{Machine Learning for Dummies}, For Dummies, 2016.;%
Mitchell, \textit{Machine Learning}, selected chapters on learning basics and evaluation.;%
Scikit-learn documentation, accessed 2026-02-28.%
}

\newcommand{\AIMLUnitIVSources}{%
Mueller and Massaron, \textit{Machine Learning for Dummies}, For Dummies, 2016.;%
James, Witten, Hastie, and Tibshirani, \textit{An Introduction to Statistical Learning} (2e), selected chapters.;%
Scikit-learn documentation, accessed 2026-02-28.%
}

\newcommand{\AIMLUnitVSources}{%
NIST, \textit{AI Risk Management Framework (AI RMF 1.0)}, 2023.;%
Barocas, Hardt, and Narayanan, \textit{Fairness and Machine Learning} (online book), selected sections.;%
Knaflic, \textit{Storytelling with Data}, selected chapters on communicating results.%
}
%
}{%
  % Faculty copies live under `private/` so their compile CWD is different.
  \IfFileExists{../../../../public/_shared/aiml-sources.tex}{%
    % Source strings for Elements of AIML (used by slides + notes).

\newcommand{\AIMLRepoURL}{https://github.com/tali7c/Elements-of-AIML}

\newcommand{\AIMLLabSources}{%
Provost and Fawcett, \textit{Data Science for Business}.;%
McKinney, \textit{Python for Data Analysis}.;%
WEKA Documentation (Waikato Environment for Knowledge Analysis), accessed 2026-02-28.;%
Matplotlib and Seaborn documentation, accessed 2026-02-28.%
}

\newcommand{\AIMLUnitISources}{%
Rich and Knight, \textit{Artificial Intelligence}, McGraw Hill, 2017.;%
Russell and Norvig, \textit{Artificial Intelligence: A Modern Approach} (4e), Ch. 1--2 (Introduction, Intelligent Agents).;%
Poole and Mackworth, \textit{Artificial Intelligence: Foundations of Computational Agents} (2e), selected sections.%
}

\newcommand{\AIMLUnitIISources}{%
Russell and Norvig, \textit{Artificial Intelligence: A Modern Approach} (4e), Ch. 7--9 (Logical Agents, First-Order Logic, Inference).;%
Huth and Ryan, \textit{Logic in Computer Science} (2e), selected sections on propositional and first-order logic.;%
Genesereth and Nilsson, \textit{Logical Foundations of Artificial Intelligence}, selected chapters.%
}

\newcommand{\AIMLUnitIIISources}{%
Mueller and Massaron, \textit{Machine Learning for Dummies}, For Dummies, 2016.;%
Mitchell, \textit{Machine Learning}, selected chapters on learning basics and evaluation.;%
Scikit-learn documentation, accessed 2026-02-28.%
}

\newcommand{\AIMLUnitIVSources}{%
Mueller and Massaron, \textit{Machine Learning for Dummies}, For Dummies, 2016.;%
James, Witten, Hastie, and Tibshirani, \textit{An Introduction to Statistical Learning} (2e), selected chapters.;%
Scikit-learn documentation, accessed 2026-02-28.%
}

\newcommand{\AIMLUnitVSources}{%
NIST, \textit{AI Risk Management Framework (AI RMF 1.0)}, 2023.;%
Barocas, Hardt, and Narayanan, \textit{Fairness and Machine Learning} (online book), selected sections.;%
Knaflic, \textit{Storytelling with Data}, selected chapters on communicating results.%
}
%
  }{%
    % Source strings for Elements of AIML (used by slides + notes).

\newcommand{\AIMLRepoURL}{https://github.com/tali7c/Elements-of-AIML}

\newcommand{\AIMLLabSources}{%
Provost and Fawcett, \textit{Data Science for Business}.;%
McKinney, \textit{Python for Data Analysis}.;%
WEKA Documentation (Waikato Environment for Knowledge Analysis), accessed 2026-02-28.;%
Matplotlib and Seaborn documentation, accessed 2026-02-28.%
}

\newcommand{\AIMLUnitISources}{%
Rich and Knight, \textit{Artificial Intelligence}, McGraw Hill, 2017.;%
Russell and Norvig, \textit{Artificial Intelligence: A Modern Approach} (4e), Ch. 1--2 (Introduction, Intelligent Agents).;%
Poole and Mackworth, \textit{Artificial Intelligence: Foundations of Computational Agents} (2e), selected sections.%
}

\newcommand{\AIMLUnitIISources}{%
Russell and Norvig, \textit{Artificial Intelligence: A Modern Approach} (4e), Ch. 7--9 (Logical Agents, First-Order Logic, Inference).;%
Huth and Ryan, \textit{Logic in Computer Science} (2e), selected sections on propositional and first-order logic.;%
Genesereth and Nilsson, \textit{Logical Foundations of Artificial Intelligence}, selected chapters.%
}

\newcommand{\AIMLUnitIIISources}{%
Mueller and Massaron, \textit{Machine Learning for Dummies}, For Dummies, 2016.;%
Mitchell, \textit{Machine Learning}, selected chapters on learning basics and evaluation.;%
Scikit-learn documentation, accessed 2026-02-28.%
}

\newcommand{\AIMLUnitIVSources}{%
Mueller and Massaron, \textit{Machine Learning for Dummies}, For Dummies, 2016.;%
James, Witten, Hastie, and Tibshirani, \textit{An Introduction to Statistical Learning} (2e), selected chapters.;%
Scikit-learn documentation, accessed 2026-02-28.%
}

\newcommand{\AIMLUnitVSources}{%
NIST, \textit{AI Risk Management Framework (AI RMF 1.0)}, 2023.;%
Barocas, Hardt, and Narayanan, \textit{Fairness and Machine Learning} (online book), selected sections.;%
Knaflic, \textit{Storytelling with Data}, selected chapters on communicating results.%
}
%
  }%
}


\title{Elements of AIML Lab\\Experiment 01 (After-submission Reference): WEKA for Regression}
\author{Tofik Ali}
\date{\today}

\begin{document}
\maketitle
\tableofcontents

\section{How to use this reference}
This document is meant to be shared \textbf{after you submit} Experiment~01.
Use it to cross-check:
\begin{itemize}[leftmargin=*]
  \item whether you selected the correct \textbf{class attribute} (numeric target),
  \item whether you evaluated correctly (CV or hold-out, not training-set),
  \item whether you interpreted MAE/RMSE in the \textbf{units of the target}.
\end{itemize}
Numbers in WEKA depend on the dataset and on random seeds/parameters, so your exact values may differ.

\section{Quick recap: what regression means}
Regression is \textbf{supervised learning} where the target is \textbf{numeric} (continuous), e.g., price, time, temperature.
In WEKA, regression is done in the \textbf{Classify} tab using a \textbf{regressor} (not a classifier).

\subsection{Key terms (minimum)}
\begin{itemize}[leftmargin=*]
  \item \textbf{Model}: a function that maps features $x$ to a numeric prediction $\hat{y}$.
  \item \textbf{MAE}: average absolute error (same units as $y$).
  \item \textbf{RMSE}: square-root of average squared error (same units as $y$; punishes large errors).
  \item \textbf{Cross-validation (CV)}: repeat train/test on folds and average to get a stable estimate.
  \item \textbf{Overfitting}: model performs well on training data but poorly on unseen data.
\end{itemize}

\section{Worked example in WEKA (end-to-end)}
\subsection{Dataset choice (example)}
If you do not have a dataset, WEKA often ships with sample datasets under a \texttt{data/} folder.
For regression, a common example is a housing dataset where the target is a numeric value.

\subsection{Step 1: load data and set class attribute}
\begin{enumerate}[leftmargin=*]
  \item Open \textbf{WEKA} $\rightarrow$ \textbf{Explorer}.
  \item \textbf{Preprocess} tab $\rightarrow$ \textbf{Open file...} $\rightarrow$ load \texttt{.arff} or \texttt{.csv}.
  \item In the right panel, set \textbf{Class} to the \textbf{target} column (the value to predict).
\end{enumerate}
\textbf{Cross-check:} the class attribute must be \textbf{numeric}. If it is nominal, you are doing classification.

\subsection{Step 2: minimal preprocessing (recommended)}
Typical safe steps:
\begin{itemize}[leftmargin=*]
  \item Missing values $\rightarrow$ \texttt{Filter} $\rightarrow$ \texttt{unsupervised/attribute/ReplaceMissingValues}.
  \item Remove obvious IDs (e.g., \texttt{CustomerID}) that do not represent a measurable feature.
\end{itemize}
Record any filter you apply (preprocessing is part of the final pipeline).

\subsection{Step 3: train and evaluate two regressors}
\paragraph{Model A: Linear Regression.}
\begin{enumerate}[leftmargin=*]
  \item Go to \textbf{Classify}.
  \item \textbf{Choose} $\rightarrow$ \texttt{functions/LinearRegression}.
  \item \textbf{Test options} $\rightarrow$ choose \textbf{Cross-validation} (10 folds) \emph{or} \textbf{Percentage split} (e.g., 70/30).
  \item Click \textbf{Start}.
\end{enumerate}
In the output, note MAE and RMSE (and correlation coefficient if shown).

\paragraph{Model B: Random Forest (a strong baseline).}
\begin{enumerate}[leftmargin=*]
  \item \textbf{Choose} $\rightarrow$ \texttt{trees/RandomForest}.
  \item Keep the \textbf{same} test option as Model~A (fair comparison).
  \item Click \textbf{Start}.
\end{enumerate}

\subsection{Step 4: interpret the metrics correctly}
Write metric interpretations in your dataset's units. Example:
\begin{itemize}[leftmargin=*]
  \item If the target is \texttt{Price (INR)} and RMSE $= 5000$, then typical errors are about \textbf{5000 INR}.
  \item If RMSE is much larger than MAE, your model may be making a few very large mistakes (outliers / hard cases).
\end{itemize}

\subsection{Step 5: optional error visualization (useful)}
WEKA can visualize prediction errors for some models.
If available:
\begin{itemize}[leftmargin=*]
  \item in the output list, right-click the result $\rightarrow$ \textbf{Visualize classifier errors}.
\end{itemize}
This helps you explain \emph{where} the model fails (certain ranges of the target, or certain feature regions).

\section{What a strong submission looks like (cross-check)}
\subsection{Screenshots}
\begin{itemize}[leftmargin=*]
  \item Preprocess tab: dataset loaded + correct \textbf{class attribute} selected.
  \item Classify tab: regressor selected + test option selected.
  \item Output: metrics (MAE/RMSE) for at least two regressors.
\end{itemize}

\subsection{Report structure (1--2 pages)}
\begin{enumerate}[leftmargin=*]
  \item Dataset: name/source, target definition, number of rows/columns.
  \item Preprocessing: missing value handling, removed columns, any transformations.
  \item Models: Model~A and Model~B with evaluation method.
  \item Results table: MAE/RMSE (+ any other metrics you used).
  \item Observations: 3--5 clear points explaining which model is better and why.
\end{enumerate}

\section{Troubleshooting (common issues and fixes)}
\begin{itemize}[leftmargin=*]
  \item \textbf{Class attribute is wrong:} re-check the \textbf{Class} selection in Preprocess.
  \item \textbf{Evaluating on training set:} in Test options, avoid \textbf{Use training set}.
  \item \textbf{Target accidentally nominal:} if the target is stored as text/categorical, convert it to numeric (or choose a proper numeric target).
  \item \textbf{Very poor error values:} check for outliers, wrong units, or missing values not handled.
  \item \textbf{Unfair comparison:} use the same evaluation method for all models.
\end{itemize}

\notesources{\AIMLLabSources}
\end{document}

